\uuid{sb3N}
\titre{Noyau/Image/injection/Surjection}
\auteur{hanani ; mzari}
\organisation{exo7}
\licence{CC BY-NC-SA 4.0}
\datecreate{2020-01-20}
\niveau{L1}
\module{Algèbre}
\chapitre{Application linéaire}
\sousChapitre{Définition}
\difficulte{3}
\theme{Noyau et image, Noyau, Image, injection, Surjection}
\exo7qcmid{lille-2.applications-lineaires.027}

\contenu{
\begin{question}
\qclassification{107.01}
\qtags{motcle=Noyau/Image/injection/Surjection}
On considère $E$ un $\Rr$-espace vectoriel et $f$ un endomorphisme de $E$ tel que : $f^2+f+Id = 0$, où $Id$ est l'application identité de $E$. Quelles sont les assertions vraies ?
\begin{answers}
\bad{$\dim \ker f = 1$}.
\good{$f$ est injective}.
\good{$f$ est bijective et $f^{-1}=f^2$}.
\good{$f$ est  bijective et $f^{-1}=-f-Id$}.
\end{answers}
\begin{explanations} Soit $x\in E$ tel que $f(x)=0$. De l'égalité  $f^2(x)+f(x)+x = 0$, on déduit que $x=0$, donc $\dim \ker f = 0$ et donc $f$ est injective.
\vskip0mm
De l'égalité  $f^2+f+Id = 0$, on déduit que $fo(-f-Id)=Id$, donc $f$ est bijective et $f^{-1}=-f-Id=f^2$.
\end{explanations}
\end{question}
}
